\documentclass[a4paper,12pt]{article}
\usepackage[utf8]{inputenc}
\usepackage[T1]{fontenc}
\usepackage{amsmath,amssymb}
\usepackage{hyperref}
\usepackage{geometry}
\geometry{margin=2.5cm}
\title{Evaluation of \(\displaystyle I = \int_0^1 \cos(x^2+e^x)\,dx\)}
\author{}
\date{}
\begin{document}
\maketitle

\section{Introduction}
We are looking to evaluate 
\[
I = \int_0^1 \cos(x^2 + e^x)\,dx .
\]
No elementary antiderivative exists, but we can obtain an exact expression in the form of a double series.

\section{Power series expansion}
For any real number $u$,
\[
\cos u = \sum_{n=0}^{\infty} \frac{(-1)^n}{(2n)!}\, u^{2n}.
\]
By setting $u = x^2 + e^x$, we have
\[
\cos(x^2+e^x) = \sum_{n=0}^{\infty} \frac{(-1)^n}{(2n)!}\, (x^2+e^x)^{2n}.
\]

\section{Binomial expansion}
For each $n$, let us expand $(x^2+e^x)^{2n}$ using the binomial theorem:
\[
(x^2+e^x)^{2n} = \sum_{k=0}^{2n} \binom{2n}{k} (x^2)^k (e^x)^{2n-k}
= \sum_{k=0}^{2n} \binom{2n}{k} x^{2k} e^{(2n-k)x}.
\]

\section{Series/integral interchange}
The series converges normally on $[0,1]$, so we can integrate term by term:
\[
I = \sum_{n=0}^{\infty} \frac{(-1)^n}{(2n)!}
\sum_{k=0}^{2n} \binom{2n}{k}
\underbrace{\int_0^1 x^{2k} e^{(2n-k)x}\,dx}_{I_{k,n}} .
\]

\section{Exact calculation of \(I_{k,n}\)}
Let $m = 2k$ and $a = 2n-k$. We must evaluate
\[
I(m,a) = \int_0^1 x^{m} e^{a x}\,dx .
\]
\begin{itemize}
\item If $a = 0$ (i.e., $k=2n$): $\displaystyle I(m,0) = \frac{1}{m+1} = \frac{1}{4n+1}$.
\item If $a \neq 0$, we integrate by parts $m$ times or use the explicit formula
\begin{align*}
I(m,a) &= \frac{m!}{(-a)^{m+1}}\Bigl(1 - e^{a}\sum_{j=0}^{m}\frac{(-a)^j}{j!}\Bigr) \\
&= \frac{(2k)!}{(k-2n)^{2k+1}}\Bigl(1 - e^{2n-k}\sum_{j=0}^{2k}\frac{(k-2n)^j}{j!}\Bigr).
\end{align*}
\end{itemize}

\section{Exact expression of the integral}
By substituting both cases into the double sum, we obtain the exact representation:
\begin{align}
I &= \sum_{n=0}^{\infty} \frac{(-1)^n}{(2n)!}
\Bigg[ \binom{2n}{2n}\frac{1}{4n+1} \nonumber \\
&\quad + \sum_{\substack{k=0\\k\neq 2n}}^{2n} \binom{2n}{k}
\frac{(2k)!}{(k-2n)^{2k+1}}
\Bigl(1 - e^{2n-k}\sum_{j=0}^{2k}\frac{(k-2n)^j}{j!}\Bigr) \Bigg]. \label{eq:exact}
\end{align}
This expression is exact and only involves factorials, binomial coefficients, and the exponential $e$.

\section{Numerical evaluation}
\subsection{Direct evaluation}
Using an adaptive quadrature method (Gauss-Kronrod), we find
\[
I \approx -0.303666963947415 \quad (\text{estimated error } \sim 10^{-15}).
\]

\subsection{Evaluation of the exact series}
The naive calculation of (\ref{eq:exact}) in double precision fails quickly: the terms $e^{2n-k}$ and the exponential sums are very large, and their subtraction causes catastrophic cancellation. For $n\ge 10$, the numerical results blow up.

Using multi-precision arithmetic (50 decimal places) and a robust numerical integration of each $I_{k,n}$ (rather than the unstable closed formula), the double series converges to:
\[
I = -0.303666963947415171254014036015\ldots
\]
The two approaches match up to machine precision, validating the exact formula.

\section{Retained approximate value}
\[
\boxed{\int_0^1 \cos(x^2+e^x)\,dx \simeq -0.303666963947415}
\]
with an uncertainty of less than $10^{-14}$.

\section*{Remark}
The expression (\ref{eq:exact}) provides an \textbf{exact value} of the integral, although it cannot be reduced to an elementary constant. For practical calculations, direct numerical integration is preferred.

\end{document}